%! arara: clean: {
%! arara: --> extensions:
%! arara: --> ['aux', 'bbl', 'bcf', 'blg', 'log', 'nav', 'out', 'pdf', 'run.xml', 'snm', 'toc']
%! arara: --> }
%! arara: lualatex: {
%! arara: --> shell: yes,
%! arara: --> draft: yes,
%! arara: --> interaction: batchmode
%! arara: --> }
%! arara: biber
% arara: lualatex: {
% arara: --> shell: yes,
% arara: --> draft: no,
% arara: --> interaction: batchmode
% arara: --> }
% arara: lualatex: {
% arara: --> shell: yes,
% arara: --> draft: no,
% arara: --> interaction: batchmode
% arara: --> }
%! arara: clean: {
%! arara: --> extensions:
%! arara: --> ['aux', 'bbl', 'bcf', 'blg', 'log', 'nav', 'out', 'run.xml', 'snm', 'toc']
%! arara: --> }
\documentclass[
	spanish,
	addpoints,
	noanswers,
	a4paper,
	8pt,
	leqno
]{exam}
\usepackage[T1]{fontenc}
\usepackage[spanish]{babel}
\usepackage{libertine}
\usepackage{geometry}
\geometry{
    left = 2cm,
    right = 2cm,
    top = 2cm,
    bottom = 3cm
}
\usepackage[intlimits]{mathtools}
\usepackage{amssymb}
\usepackage{diffcoeff}
\usepackage{amsthm}
\usepackage[svgnames]{xcolor}
\usepackage{siunitx}
\usepackage{tikz}
\usepackage[shortlabels]{enumitem}
\usepackage{multicol}
\usepackage{hyperref}
\theoremstyle{definition}
\newtheorem{definition}{Definición}
\pagestyle{headandfoot}

\usepackage{minted}
\renewcommand{\solutiontitle}{\noindent\textbf{Solución}\par\noindent}
\renewcommand\listingscaption{Listado}
\newcommand{\unmarkedfntext}[1]{%
    \begingroup
    \renewcommand\thefootnote{}\footnote{#1}%
    \addtocounter{footnote}{-1}%
    \endgroup
}

\difdef{c}{L}{op-symbol=\mathop{}\!\mathbin\bigtriangleup}

\pointpoints{Punto}{Puntos}
\bracketedpoints

\difdef {f,fp,s,sp} {|}{
    outer-Ldelim =\left.,
    outer-Rdelim=\right|,
    sub-nudge=0mu
}
\footer{}{Página \thepage\ de \numpages}{}


\pointpoints{Punto}{Puntos}
\bracketedpoints
\begin{document}
\begin{center}
	\sffamily\bfseries\Large
	Métodos Numéricos para Ecuaciones Diferenciales Parciales~CM 032 A\\
	Laboratorio N$^{\circ}1$
\end{center}

\begin{center}
	\fbox{\fbox{\parbox{5.5in}
			{\centering
				Responda las preguntas en los espacios provistos en las hojas
				de preguntas.
				Los argumentos y la claridad de las respuestas se
				considerarán en la puntuación final.
			}}}
\end{center}
\vspace{0.1in}
\makebox[\textwidth]{Nombres y apellidos:\enspace\hrulefill}
\vspace{0.2in}
\makebox[\textwidth]{Nombres y apellidos del profesor:\enspace Fidel Jara Huanca.\hfill}

\begin{questions}
	% \header{Math 115}{Second Exam}{July 4, 1776}
	\question\label{question:1}

	Un enfoque para derivar fórmulas de diferencias para aproximar
	derivadas es el método de coeficientes indeterminados.
	Supongamos que $f$ es una función suave en $\mathbb{R}$.
	Sea $h>0$.
	Determine los coeficientes $a$, $b$ y $c$ de modo que
	\begin{equation*}
		af\left(x+h\right)+
		bf\left(x\right)+
		cf\left(x-h\right)
	\end{equation*}
	es una aproximación de $f^{\prime}\left(x\right)$ con un orden como
	lo más alto posible; es decir, elija $a$, $b$ y $c$ tales que
	\begin{equation*}
		\left|
		af\left(x+h\right)+bf\left(x\right)+
		cf\left(x-h\right)-f^{\prime}\left(x\right)
		\right|\leq
		\mathcal{O}\left(h^{p}\right).
	\end{equation*}
	% https://www.colorado.edu/amath/sites/default/files/attached-files/mathcomp_88_fd_formulas.pdf
	% https://www.phys.uconn.edu/~rozman/Courses/m3510_20f/downloads/fornberg.pdf
	% https://www.colorado.edu/amath/sites/default/files/attached-files/sirev_cl.pdf
	% https://www.phys.uconn.edu/~rozman/Courses/m3510_20f/downloads/hw08.pdf
	% https://www.youtube.com/watch?v=M_ixanyL7IA
	% https://en.wikipedia.org/wiki/Numerical_differentiation#Complex-variable_methods

	\question

	Haz el mismo problema del Ejercicio~\ref{question:1} con
	$f^{\prime}\left(x\right)$ reemplazado por
	$f^{\prime\prime}\left(x\right)$.

	\question

	¿Es posible utilizar
	\begin{equation*}
		af\left(x+h\right)+
		bf\left(x\right)+
		cf\left(x-h\right)
	\end{equation*}
	con coeficientes elegidos adecuadamente para aproximar
	$f^{\prime\prime\prime}\left(x\right)$?
	¿Cuántos valores de función se necesitan para aproximar
	$f^{\prime\prime\prime}\left(x\right)$?

	% \begin{definition}{\em Árbol}
	% Un árbol $T$.
	% \end{definition}

	\question\label{question:4}

	Para el problema del valor inicial de la ecuación de onda
	unidireccional
	\begin{equation}\label{eq:IVP-one-way-PDE}
		\begin{cases}\displaystyle
			\diffp{u}{t}=
			a\diffp{u}{x}+
			f   &
			\text{en }\mathbb{R}\times\left(0,\infty\right). \\
			u=f &
			\text{en }\mathbb{R}\times\left\{t=0\right\}.
		\end{cases}
	\end{equation}
	Donde $a\in\mathbb{R}$ es una constante, deriva alguna diferencia
	esquemas basados en varias combinaciones de aproximaciones en
	diferencias de la derivada temporal y la derivada espacial.

	\question

	La idea del esquema de Lax-Wendroff para resolver el problema del
	valor inicial~\eqref{eq:IVP-one-way-PDE} del
	Ejercicio~\ref{question:4} es el siguiente.
	Comience con la expansión de Taylor
	\begin{equation}\label{eq:taylor}
		u\left(x_{j},t_{m+1}\right)\approx
		u\left(x_{j},t_{m}\right)+
		h_{t}\diffp{u}{t}\left(x_{j},t_{m}\right)+
		\frac{h^{2}_{t}}{2}
		\diffp[2]{u}{t}\left(x_{j},t_{m}\right).
	\end{equation}
	De la ecuación diferencial tenemos
	\begin{equation*}
		\diffp{u}{t}=
		-a\diffp{u}{x}+
		f
	\end{equation*}
	y
	\begin{equation*}
		\diffp[2]{u}{t}=
		a^{2}\diffp[2]{u}{x}-
		a\diffp{f}{x}+
		\diffp{f}{t}.
	\end{equation*}
	Utilice estas relaciones para reemplazar las derivadas del tiempo
	en el lado derecho de~\eqref{eq:taylor}.
	Luego reemplace la primera y segunda derivadas espaciales por
	diferencias centrales.
	Finalmente, reemplace $\diffp{f}{x}$ por una diferencia central y
	$\diffp{f}{t}$ por una diferencia directa.
	Siga las instrucciones anteriores para derivar el esquema de
	Lax-Wendroff para resolver~\eqref{eq:IVP-one-way-PDE}.

	\question

	Dar esquemas hacia adelante, hacia atrás y Crank-Nicolson para el problema de valor inicial frontera
	\begin{equation*}
		\begin{cases}\displaystyle
			\diffp{u}{t}=
			\nu_{1}\diffp[2]{u}{x}+
			\nu_{2}\diffp[2]{u}{y}+
			f\left(x,y,t\right) &
			\text{en }\left(0,a\right)\times\left(0,b\right)\times\left(0,T\right). \\
			u=g_{l}             &
			\text{en }\left\{0\right\}\times\left[0,b\right]\times\left[0,T\right]. \\
			u=g_{r}             &
			\text{en }\left\{a\right\}\times\left[0,b\right]\times\left[0,T\right]. \\
			u=g_{b}             &
			\text{en }\left[0,a\right]\times\left\{0\right\}\times\left[0,T\right]. \\
			u=g_{t}             &
			\text{en }\left[0,a\right]\times\left\{b\right\}\times\left[0,T\right]. \\
			u=u_{0}             &
			\text{en }\left[0,a\right]\times\left[0,b\right]\times\left\{t=0\right\}.
		\end{cases}
	\end{equation*}
	Aquí, para los datos dados, $\nu_{1}$, $\nu_{2}$, $a$, $b$, $T$ son
	números positivos, y $f$, $g_{l}$, $g_{r}$, $g_{b}$, $g_{t}$,
	$u_{0}$ son funciones continuas de sus argumentos.

	\question

	Considere el problema del valor inicial frontera de una ecuación hiperbólica
	\begin{equation*}
		\begin{cases}\displaystyle
			\diffp[2]{u}{t}=
			\nu\diffp[2]{u}{x}+
			f\left(x,t\right)  &
			\text{en }\left(0,a\right)\times\left(0,T\right).   \\
			u=g_{1}            &
			\text{en }\left\{0\right\}\times\left[0,T\right].   \\
			u=g_{2}            &
			\text{en }\left\{a\right\}\times\left[0,T\right].   \\
			u=u_{0}            &
			\text{en }\left[0,a\right]\times\left\{t=0\right\}. \\
			\diffp{u}{t}=v_{0} &
			\text{en }\left[0,a\right]\times\left\{t=0\right\}.
		\end{cases}
	\end{equation*}
	Aquí, para los datos dados, $\nu$, $a$, $T$ son números positivos y
	$f$, $g_{1}$, $g_{2}$, $u_{0}$, $v_{0}$ son funciones continuas de sus argumentos.
	Analice cómo formar un esquema de diferencias finitas razonable para resolver el problema.
\end{questions}
\vfill
\begin{flushright}\bfseries
	Facultad de Ciencias\\[2mm]
	Universidad Nacional de Ingeniería\\[2mm]
	\today%\unmarkedfntext{Cada pregunta vale 5 puntos.}
\end{flushright}
\end{document}

